\section{Introduction}
\label{sec:intro}

State constitutional redistricting criteria are typically listed in priority
order.  For congressional maps, the hierarchy is: (1) equal population
(Wesberry v.\ Sanders~\citep{wesberry1964}); (2) VRA compliance
(\emph{Thornburg v.\ Gingles}~\citep{gingles1986}); (3) contiguity;
(4) compactness; and (5) preservation of political subdivisions (counties
and municipalities).

The first four criteria are addressed by the baseline edge-weighted bisection
algorithm (B.2).  The fifth --- \emph{county preservation} --- requires
explicit treatment.  Most states have constitutional or statutory provisions
requiring maps to ``preserve political subdivisions to the extent possible,''
a phrase that appears in some form in 34 state constitutions~\citep{ncsl2021}.

\subsection*{Why existing approaches fail}

Two naive approaches to county preservation fail in practice:

\paragraph{Hard constraints.}
Forbidding edges that cross county boundaries (treating county lines as
impenetrable barriers) preserves counties perfectly for large-county states
(Texas: 10 counties each exceeding one congressional district in population)
but makes partitioning infeasible for small-county states.
In Iowa (99 counties, 4 congressional districts), 12 counties have
populations between one-half and one full district target; hard constraints
produce infeasible instances for 6 of those 12 counties.

\paragraph{No treatment.}
Ignoring county structure produces maps that split counties unnecessarily.
The B.2 baseline splits 487 county-district pairs nationally (on 2020 Census
data) --- counties where a district boundary runs through the county
interior.  Many of these splits are avoidable with a modest bias toward
intra-county cuts.

\subsection*{Proposed approach}

We introduce \emph{county-sticky} edge weighting: edges between census
tracts within the same county receive a multiplier $\alpha_c \geq 1$,
making intra-county cuts more expensive than inter-county cuts.
This soft-constraint approach avoids infeasibility while biasing the
algorithm toward county-respecting partitions.

\subsection*{Contributions}

This paper makes four contributions:

\begin{enumerate}
  \item We formalise the county-sticky weighting scheme and prove that it
    preserves the planarity and connectivity properties required for METIS.
  \item We sweep $\alpha_c \in \{1.0, 1.5, 2.0, 3.0, 5.0, 10.0\}$ across
    all 50 states and identify $\alpha_c = 3.0$ as the Pareto-optimal
    choice (minimising county splits subject to minimal compactness loss).
  \item We provide state-level county-split reduction results, with detailed
    case studies for Iowa ($k{=}4$) and Texas ($k{=}38$).
  \item We connect the results to the DIA's algorithm specification
    (\texttt{--weights-override county}) and state constitutional requirements.
\end{enumerate}

\subsection*{Paper structure}

Section~\ref{sec:related} reviews subdivision preservation in the
redistricting literature.
Section~\ref{sec:methodology} formalises the county-sticky formulation
and the parameter sweep design.
Section~\ref{sec:results} presents the $\alpha_c$ sweep and state-level
results.
Section~\ref{sec:discussion} connects the results to legal requirements.
Section~\ref{sec:conclusion} concludes.
